\begin{theorem}\label{night:P3K.6}
\StatementRecord{night:P3K.6}{false as stated; correction proved}{S6\_06\_BLIND: GPT-6 Astra (OpenAI, ChatGPT Pro mode, effort Pro), recorded re-derivation in the checkpoint dated 18 November 2026}{\cref{night:P3F.10}, \cref{night:P3F.13}, \cref{night:P3F.9}}\par\noindent\textbf{Physical-line correction, 17 November 2026.} The meridian-one Thom line is the Hermitian dual of the physical line. Its page restriction is \(-c_1(M)\). The embedded trace, collars and integer uniqueness are retained; the positive page-line assertion is refuted by the explicit dual arrow in K68. The following earlier assertion and its argument are retained as the located historical claim; their corrected current consequences are proved in K68--69.\medskip\par 
Let one positive gamma turn denote the actual parameter cover whose
angular clutch is P3F.13.  The orbit family of this cover is the product
\(Y\times S^1_\theta\), with \(Y=S^5\), and its fixed trace is

\[
 T_F=\{(R_\theta(p),\theta):p\in F\},
\]

where \(R_\theta\) is the gamma rotation on the cusp orbit product of
P3F.9.  Use its product trivialization preserving the Seifert gamma
coordinate.  In these coordinates \(R_\theta\) is pure gamma
rotation, including the explicitly trivialized source-translation
suspension.  There is a compact properly embedded oriented four-manifold
\(X\subset Z=Y\times D^2\), with boundary \(T_F\), assembled from the
actual page \(V\) and a product compression trace.  Its complement
has a unique integral class \(x\) with meridian value one.  On the
outer boundary this class is the actual delta-circle Euler class.
On the graph copy of \(V\), its restriction to the framed normal
section is exactly \(c_1(M)\).
\end{theorem}

\begin{proof}
\begin{enumerate}[label=\arabic*.]
\item \par\medskip\noindent\textbf{The actual orbit trivialization.}\enspace The ordinary and finite
orbit charts are independent of the parameter angle: the variation
of their period and affine-twist columns is in the real delta
direction.  The exact subtraction in P3F.13 leaves, on the cusp
overlap, the gamma translation \(R_\theta\); the residual delta term
is a periodic unitary gauge.  Apply the identity to the exterior
charts and \(R_\theta\) to the cusp chart.  Their overlap maps now
agree pointwise.  Since \(R_1=R_0\), these identifications give a
periodic trivialization of the orbit family.  The fixed sphere in
the cusp chart becomes precisely \(R_\theta F\).  This verifies the
product and the displayed trace for the actual angular family.

\par\medskip\noindent\textbf{The compression.}\enspace Choose the binding at a regular value
\(q=\epsilon>0\) in a small cusp disc.  First write the local
regular-binding Milnor function as
\(u(q)(q-\epsilon)e^{2\pi i12\gamma_{\rm old}}\), with \(u\ne0\),
choose its unit argument, and set
\(\gamma=\gamma_{\rm old}+\arg u(q)/(2\pi\,12)\).
This is an explicit base-dependent circle coordinate change,
commuting with rotation and preserving the integral gamma marking.
Make this unit correction before choosing the short arc and the
handle below.  At its node the source has
coordinates \(xy=q\).  The two normalization discs of the fixed
sphere project to \(x=0\) and \(y=0\) and are separated in the gamma
coordinate.  Keep their values and the rest of the fixed sphere
inside an arc \(J\) of length less than \(1/24\), as allowed by the
relative product construction P3F.9.  Write their actual real gamma
functions as \(g_x(z)\) and \(g_y(z)\), without changing the orbit
map or its page.

\item Use the following exact real handle in the node chart:

\[
 (z,\phi)\longmapsto
 \bigl(x=z\cos\phi,\quad y=\bar z\sin\phi,\quad
       \gamma=g(z,\phi)\bigr),
 \qquad 0\leq\phi\leq\pi/2,
\]

\item Here \(g\) interpolates from \(g_x(z)\) to \(g_y(\bar z)\), with
values in \(J\); for example use their convex linear interpolation
in \(2\phi/\pi\).  The two center values are distinct, so
\(g(0,\phi)\) is strictly monotone.  For \(z\ne0\), the pair
\((x,y)\) determines \((z,\phi)\) by its moduli and phases.  At
\(z=0\), gamma determines \(\phi\), and at least one of
\(\sin\phi,\cos\phi\) is nonzero.  These facts also give full
rank to the derivative, so the map is an embedding at the core.
Its exact base value is

\[
 q=xy=|z|^2\cos\phi\sin\phi\geq0.
\]

\item Restrict to \(|x|,|y|\leq a\) and \(q\leq\epsilon\).
The faces \(\phi=0,\pi/2\) are the two normalization discs at
separated gamma values.  The face \(q=\epsilon\) is the smoothing
annulus, and the remaining faces \(|x|=a\), \(|y|=a\) match the
transported exterior of the source family along the real path.
Their gamma functions extend over that exterior with values in
\(J\), since \(J\) is an interval.  Rounding only these product
corners gives the oriented compression body.  Its two disc-face
orientations are the complex normalization orientations: the
conjugation on the second face is canceled by its opposite boundary
orientation.  Equivalently, reversing it attaches the positive
smoothing one-handle to \(F\).  Its regular boundary is the source
torus embedded as a graph \(\gamma=g_\epsilon(p)\), with the
globally defined real-valued function \(g_\epsilon\) taking values
in \(J\).  Choose a collar-preserving height function to obtain
an embedded three-dimensional compression cobordism

\[
 W\hookrightarrow Y\times[r_0,1],\qquad
 \partial W=F\sqcup(-E_\epsilon).
\]

\item Its projection lies over \(0\leq q\leq\epsilon\) and inside \(J\),
by the exact displayed formula.  The trace is used in the product
with its height interval.  Topologically \(W\) is a solid torus
with an open three-ball removed.  Its meridian on the incoming
torus is the source vanishing cycle \(a\).  The explicit handle
also fixes the framing and avoids an uncontrolled normal
perturbation of the base coordinate.

\par\medskip\noindent\textbf{Choose the page and embed the graph.}\enspace The weighted circle
has degree twelve on the regular-binding Milnor function.  In the
already chosen unit-corrected coordinate, the page equation is

\[
 \arg(q-\epsilon)+2\pi\,12\gamma=\varphi.
\]

\item Along the real path \(0\leq q<\epsilon\), its twelve gamma values
are constant.  Choose \(\varphi\) so that all these values miss
\(J\), and so that the short arc \(12J\) lies strictly on one side
of the page near the binding.  This is possible because of the
chosen length of \(J\).  The interior of the actual page \(V\)
then misses the entire projected compression, not just \(F\).
At the binding the initial torus meets \(\partial V\); pushing the
page a small distance to the opposite page-normal side removes this
intersection as well.  Denote this pushoff by \(V^+\).  It is
isotopic to \(V\) in \(Y\setminus F\) and is disjoint from the
projected compression.

\item On \(\partial V=E_\epsilon\times S^1_\gamma\), define
\(f(p,\gamma)=r_0e^{2\pi i(\gamma-g_\epsilon(p))}\).
This is the exact inverse angular coordinate of the rotating
incoming graph.  Its integral marking is the gamma marking,
because \(g_\epsilon\) is a globally defined real-valued function.
Extend this to a smooth map
\(f:V\to D^2_{r_0}\): on a boundary collar multiply the displayed
value by a cutoff that decreases to zero and is flat at zero,
then set \(f=0\) on the remaining interior.  Its graph in
\(Y\times D^2\) is embedded because its projection to \(Y\) is the
original embedded page.  In the outer annulus attach

\[
 W\times S^1\longrightarrow Y\times D^2,
 \qquad ((y,r),\theta)\longmapsto
 (R_\theta y,r e^{2\pi i\theta}).
\]

\item The polar coordinate records both \(r\) and \(\theta\), so this
map is an embedding.  Its incoming boundary is exactly the graph
of \(f|_{\partial V}\), and its outgoing boundary is \(T_F\).
Choose the radial derivative of the graph collar nonzero.  Extend
that collar outward and adjust only the outer-annulus collar of
the compression to meet it.  In a normal chart of the common
binding these are graphs over the same radial collar, so smooth
interpolation gives an embedded union.  All this adjustment is
arbitrarily small and leaves the graph inside the closed inner
disc unchanged.  At \(\theta=0\) it is chosen on the compression
side of the page; the fixed opposite pushoff \(V^+\) remains
disjoint.  This is the asserted \(X\).  Novikov additivity gives
\(\sigma(X)=\sigma(V)\): the added product \(W\times S^1\) has zero
signature, as is also seen by reversing its circle factor.

\par\medskip\noindent\textbf{The euler line, including its clutch.}\enspace The integral Thom
sequence for the proper codimension-three embedding \(X\subset Z\)
contains

\[
 H^2(Z;\mathbf Z)=0\longrightarrow H^2(Z\setminus X;\mathbf Z)
 \longrightarrow H^0(X;\mathbf Z)=\mathbf Z
 \longrightarrow H^3(Z;\mathbf Z)=0.
\]

\item It gives exactly one class \(x\) with the chosen meridian value one.
The identical calculation for
\(T_F\subset Y\times S^1\) identifies
\(H^2((Y\times S^1)\setminus T_F;\mathbf Z)=\mathbf Z\)
by the same meridian.  The actual delta bundle has that value,
so its first Chern class is the restriction of \(x\).  Classification
of principal circle bundles by integral first Chern class supplies
an isomorphism of these specified bundles.  This is an integral
identification on the whole complement; it retains the angular
clutch of degree minus twelve, rather than deleting its gauge.

\item The normal bundle of the graph copy of \(V\) is the sum of its
oriented page-normal real line and the real tangent plane of the
parameter disc.  The graph derivative gives an upper-triangular
identification with this sum, hence a framing.  Take the section
pointing toward \(V^+\).  Its pushoff first moves its disc coordinate
to zero inside the radius \(r_0\); its \(Y\)-coordinate is off the
graph throughout.  It then moves radially from zero to the outer
parameter point \(\theta=0\).  In the outer annulus the trace at
\(\theta=0\) is exactly \(W\), whose projection misses \(V^+\).
This is a homotopy in \(Z\setminus X\) to the pushoff of the actual
page in the original parameter-zero exterior.  The restriction of
\(x\) along that normal section is therefore

\[
 s^*x=c_1(L_\delta|_V)=c_1(M).
\]

\item The last equality is the exact circle-line comparison P3F.10.
Boundary collar choices change only its boundary gauge, and their
relative square is zero by the relative cup-product formula.
This completes the construction and the meridian comparison.

\end{enumerate}
\end{proof}